\documentclass[10pt,a4paper]{amsart}
\usepackage{amsmath,amssymb,amsthm,graphicx}
\usepackage{verbatim}
\newenvironment{reference}{\comment}{\endcomment}
%\newenvironment{reference}{}{}
\newenvironment{slogan}{\comment}{\endcomment}
\newenvironment{history}{\comment}{\endcomment}

% For commutative diagrams we use Xy-pic
\usepackage[all]{xy}

% We use 2cell for 2-commutative diagrams.
\xyoption{2cell}
% Theorem environments.
%
\theoremstyle{plain}
\newtheorem{theorem}[subsection]{Theorem}
\newtheorem{proposition}[subsection]{Proposition}
\newtheorem{lemma}[subsection]{Lemma}

\theoremstyle{definition}
\newtheorem{definition}[subsection]{Definition}
\newtheorem{example}[subsection]{Example}
\newtheorem{exercise}[subsection]{Exercise}
\newtheorem{situation}[subsection]{Situation}

\theoremstyle{remark}
\newtheorem{remark}[subsection]{Remark}
\newtheorem{remarks}[subsection]{Remarks}

\numberwithin{equation}{subsection}

% Macros
%
\def\lim{\mathop{\mathrm{lim}}\nolimits}
\def\colim{\mathop{\mathrm{colim}}\nolimits}
\def\Spec{\mathop{\mathrm{Spec}}}
\def\Hom{\mathop{\mathrm{Hom}}\nolimits}
\def\Ext{\mathop{\mathrm{Ext}}\nolimits}
\def\SheafHom{\mathop{\mathcal{H}\!\mathit{om}}\nolimits}
\def\SheafExt{\mathop{\mathcal{E}\!\mathit{xt}}\nolimits}
\def\Sch{\mathit{Sch}}
\def\Mor{\mathop{\mathrm{Mor}}\nolimits}
\def\Ob{\mathop{\mathrm{Ob}}\nolimits}
\def\Sh{\mathop{\mathit{Sh}}\nolimits}
\def\NL{\mathop{N\!L}\nolimits}
\def\CH{\mathop{\mathrm{CH}}\nolimits}
\def\proetale{{pro\text{-}\acute{e}tale}}
\def\etale{{\acute{e}tale}}
\def\QCoh{\mathit{QCoh}}
\def\Ker{\mathop{\mathrm{Ker}}}
\def\Im{\mathop{\mathrm{Im}}}
\def\Coker{\mathop{\mathrm{Coker}}}
\def\Coim{\mathop{\mathrm{Coim}}}

% Boxtimes
%
\DeclareMathSymbol{\boxtimes}{\mathbin}{AMSa}{"02}

%
% Macros for moduli stacks/spaces
%
\def\QCohstack{\mathcal{QC}\!\mathit{oh}}
\def\Cohstack{\mathcal{C}\!\mathit{oh}}
\def\Spacesstack{\mathcal{S}\!\mathit{paces}}
\def\Quotfunctor{\mathrm{Quot}}
\def\Hilbfunctor{\mathrm{Hilb}}
\def\Curvesstack{\mathcal{C}\!\mathit{urves}}
\def\Polarizedstack{\mathcal{P}\!\mathit{olarized}}
\def\Complexesstack{\mathcal{C}\!\mathit{omplexes}}
% \Pic is the operator that assigns to X its picard group, usage \Pic(X)
% \Picardstack_{X/B} denotes the Picard stack of X over B
% \Picardfunctor_{X/B} denotes the Picard functor of X over B
\def\Pic{\mathop{\mathrm{Pic}}\nolimits}
\def\Picardstack{\mathcal{P}\!\mathit{ic}}
\def\Picardfunctor{\mathrm{Pic}}
\def\Deformationcategory{\mathcal{D}\!\mathit{ef}}
\usepackage{cleveref}
\crefname{originallemma}{original lemma}{original lemmas}
\Crefname{originallemma}{Original lemma}{Original lemmas}
\makeatletter
\newlabel{lemma-base-change-affine}{{03WI}{source}{}{sourcepointer-1}{}}
\newlabel{lemma-base-change-affine@cref}{{[originallemma][0][]03WI}{[1][0][]source}{}{}{}}
\makeatother
\crefname{originallemma}{original lemma}{original lemmas}
\Crefname{originallemma}{Original lemma}{Original lemmas}
\makeatletter
\newlabel{lemma-closed-immersion-affine}{{07U2}{source}{}{sourcepointer-2}{}}
\newlabel{lemma-closed-immersion-affine@cref}{{[originallemma][0][]07U2}{[1][0][]source}{}{}{}}
\makeatother
\crefname{originallemma}{original lemma}{original lemmas}
\Crefname{originallemma}{Original lemma}{Original lemmas}
\makeatletter
\newlabel{lemma-composition-affine}{{03WH}{source}{}{sourcepointer-3}{}}
\newlabel{lemma-composition-affine@cref}{{[originallemma][0][]03WH}{[1][0][]source}{}{}{}}
\makeatother
\crefname{originaldefinition}{original definition}{original definitions}
\Crefname{originaldefinition}{Original definition}{Original definitions}
\makeatletter
\newlabel{spaces-properties-definition-separated}{{03BS}{source}{}{sourcepointer-4}{}}
\newlabel{spaces-properties-definition-separated@cref}{{[originaldefinition][0][]03BS}{[1][0][]source}{}{}{}}
\makeatother
\title{A morphism from an affine space over the base to a space with affine diagonal is affine}
\author{The Stacks Project Authors}
\date{stacks-project:spaces-morphisms.tex:local-snapshot-20261012; GNU FDL 1.2 or later}
\begin{document}
\maketitle
\setcounter{subsection}{1}
\noindent\textit{Editorial note.} All delivered mathematical text is the author's own, copied from the pinned Stacks Project chapter file spaces-morphisms.tex; no statement or proof step was rewritten, reordered or removed.\par\medskip
\noindent\textit{Editorial note.} The proof labels its steps 'Proof of (1)', 'Part (2) follows from (1)' and so on; all four assertions are handled inside the single counted proof environment.\par\medskip
\noindent\textit{Editorial note.} The four references of the counted spans are recorded as non-load-bearing navigation pointers; the argument that consumes them stays inside the counted proof span.\par\medskip
\noindent\textit{Editorial note.} Preamble: the three native excerpt spans of sources/stacks-spaces-morphisms/preamble.tex (lines 8-13, 14-18 and 154-217). The chapter file itself carries no preamble text.\par\medskip
\noindent\textit{Allowed standard prerequisites (editorial).}\begin{itemize}
\item Affineness is stable under base change. Lemma lemma-base-change-affine, used at proof lines 4133-4138. Makes both X x\_B Y -> Y and the graph (1, g) affine.
\item A composite of affine morphisms is affine. Lemma lemma-composition-affine, used at proof lines 4139-4140. Concludes affineness of g from affineness of its two factors.
\item A closed immersion is affine. Lemma lemma-closed-immersion-affine, used at proof lines 4141-4142. Turns the separatedness hypothesis into the affine diagonal hypothesis.
\item Separatedness of an algebraic space over B means the diagonal is a closed immersion. Properties of Spaces, Definition spaces-properties-definition-separated, used at proof lines 4124-4126. Identifies the two hypotheses in parts (2) and (4).
\end{itemize}\medskip
\ifdefined\hypertarget\hypertarget{sourcepointer-1}{}\fi
\noindent\textit{Source reference.} Original lemma 03WI: focused-seeds/AGP/sources/stacks-snapshot-20261011/tags.txt, source lines 5113--5113. Navigation pointer into the Stacks Project; the standard fact it supplies is applied inside the counted proof span and the argument resting on it stays there.\par\medskip
\ifdefined\hypertarget\hypertarget{sourcepointer-2}{}\fi
\noindent\textit{Source reference.} Original lemma 07U2: focused-seeds/AGP/sources/stacks-snapshot-20261011/tags.txt, source lines 9957--9957. Navigation pointer into the Stacks Project; the standard fact it supplies is applied inside the counted proof span and the argument resting on it stays there.\par\medskip
\ifdefined\hypertarget\hypertarget{sourcepointer-3}{}\fi
\noindent\textit{Source reference.} Original lemma 03WH: focused-seeds/AGP/sources/stacks-snapshot-20261011/tags.txt, source lines 5112--5112. Navigation pointer into the Stacks Project; the standard fact it supplies is applied inside the counted proof span and the argument resting on it stays there.\par\medskip
\ifdefined\hypertarget\hypertarget{sourcepointer-4}{}\fi
\noindent\textit{Source reference.} Original definition 03BS: focused-seeds/AGP/sources/stacks-snapshot-20261011/tags.txt, source lines 4359--4359. Navigation pointer into the Stacks Project; the standard fact it supplies is applied inside the counted proof span and the argument resting on it stays there.\par\medskip
\begin{lemma}
\label{lemma-affine-permanence}
Let $S$ be a scheme. Let $B$ be an algebraic space over $S$.
Suppose $g : X \to Y$ is a morphism of algebraic spaces over $B$.
\begin{enumerate}
\item If $X$ is affine over $B$ and $\Delta : Y \to Y \times_B Y$ is affine,
then $g$ is affine.
\item If $X$ is affine over $B$ and $Y$ is separated over $B$,
then $g$ is affine.
\item A morphism from an affine scheme to an algebraic space with affine
diagonal over $\mathbf{Z}$ (as in Properties of Spaces, Definition
\ref{spaces-properties-definition-separated}) is affine.
\item A morphism from an affine scheme to a separated algebraic space
is affine.
\end{enumerate}
\end{lemma}

\begin{proof}
Proof of (1). The base change $X \times_B Y \to Y$ is affine by
Lemma \ref{lemma-base-change-affine}.
The morphism $(1, g) : X \to X \times_B Y$ is the base change of
$Y \to Y \times_B Y$ by the morphism $X \times_B Y \to Y \times_B Y$.
Hence it is affine by
Lemma \ref{lemma-base-change-affine}.
The composition of affine morphisms is affine
(see Lemma \ref{lemma-composition-affine}) and (1) follows.
Part (2) follows from (1) as a closed immersion is affine
(see Lemma \ref{lemma-closed-immersion-affine}) and $Y/B$ separated
means $\Delta$ is a closed immersion. Parts (3) and (4) are special
cases of (1) and (2).
\end{proof}

\end{document}
